\begin{theorem}[有限缺陷口径下的尖点定位与离线偏移谱读出]\label{thm:xi-integrated-defect-cusp-reconstruction}
\leanverified{paper\_xi\_integrated\_defect\_cusp\_reconstruction}
作有限缺陷假设：离临界线零点多重集 \(\mathcal{Z}^-\)（定理 \ref{thm:xi-comoving-jensen-defect-kernel-decomposition}）为有限集。
令
\[
\Delta:=\{\delta(\rho):\ \rho\in\mathcal{Z}^-\}\subset(0,1),
\qquad
m(\delta):=\#\{\rho\in\mathcal{Z}^-:\ \delta(\rho)=\delta\}\in\NN\qquad(\delta\in\Delta),
\]
并对每个 \(\delta\in\Delta\) 定义阈值半径 \(\varrho_\delta:=\frac{1-\delta}{1+\delta}\)。
则对总平均缺陷 \(\mathcal{I}_D(\varrho)=\int_{\RR}D_{\zeta,x_0}(\varrho)\,\dd x_0\)（\eqref{eq:xi-integrated-defect-I}）有：
\begin{enumerate}
  \item \(\mathcal{I}_D\in C^1(0,1)\)；
  \item \(\mathcal{I}_D\) 在 \(\varrho\in(0,1)\) 处不属于 \(C^2\) 当且仅当 \(\varrho=\varrho_\delta\) 对某个 \(\delta\in\Delta\) 成立；
  \item 若 \(\varrho_\star=\varrho_{\delta_\star}\) 为上述尖点之一，则当 \(\varrho\downarrow \varrho_\star^+\) 时存在一侧展开
  \begin{equation}\label{eq:xi-integrated-defect-cusp-expansion}
  \boxed{\;
  \mathcal{I}_D(\varrho)
  =\mathcal{I}_D(\varrho_\star)+\mathcal{I}_D'(\varrho_\star)(\varrho-\varrho_\star)
  +m(\delta_\star)\,K(\delta_\star)\,(\varrho-\varrho_\star)^{3/2}
  +o\!\big((\varrho-\varrho_\star)^{3/2}\big)\;}
  \end{equation}
  其中 \(K(\delta)\) 为命题 \ref{prop:xi-single-defect-threshold-phase-transition} 的显式常数。
  因而由尖点位置与系数即可唯一读出
  \[
  \boxed{\;
  \delta_\star=\frac{1-\varrho_\star}{1+\varrho_\star},
  \qquad
  m(\delta_\star)=\frac{A_\star}{K(\delta_\star)}\in\NN\;
  }
  \]
  其中
  \(
  A_\star
  :=\lim_{\varrho\downarrow \varrho_\star^+}
  \frac{\mathcal{I}_D(\varrho)-\mathcal{I}_D(\varrho_\star)-\mathcal{I}_D'(\varrho_\star)(\varrho-\varrho_\star)}{(\varrho-\varrho_\star)^{3/2}}
  \)
  为右侧 \(3/2\) 项系数。
\end{enumerate}
\end{theorem}
\begin{proof}
由定理 \ref{thm:xi-integrated-defect-sumrule-endpoint-flux} 的有限和形式，
\[
\mathcal{I}_D(\varrho)=\sum_{\rho\in\mathcal{Z}^-}J_{\varrho}(\delta(\rho))
=\sum_{\delta\in\Delta} m(\delta)\,J_{\varrho}(\delta).
\]
对固定 \(\delta\in(0,1)\)，由命题 \ref{prop:xi-single-defect-integrated-closed-form} 的闭式可知
\(\varrho\mapsto J_{\varrho}(\delta)\)
在 \((0,1)\setminus\{\varrho_\delta\}\) 上为实解析函数，且由命题 \ref{prop:xi-single-defect-threshold-phase-transition} 得其在阈值处属于 \(C^1\) 但不属于 \(C^2\)。
故有限求和给出 \(\mathcal{I}_D\in C^1(0,1)\)，得到 (1)。
\par
再证 (2)。若 \(\varrho\neq \varrho_\delta\) 对所有 \(\delta\in\Delta\) 成立，则每一项 \(J_{\varrho}(\delta)\) 在该点邻域内皆为实解析（或恒为零），有限求和仍为 \(C^\infty\)，从而属于 \(C^2\)。
反之若 \(\varrho=\varrho_{\delta_\star}\) 对某个 \(\delta_\star\in\Delta\) 成立，则项 \(m(\delta_\star)J_{\varrho}(\delta_\star)\) 在该点右侧的二阶导按命题 \ref{prop:xi-single-defect-threshold-phase-transition} 发散，而其余各项在该点邻域内皆为 \(C^\infty\) 函数，故 \(\mathcal{I}_D\notin C^2\)。
\par
最后证 (3)。在 \(\varrho_\star=\varrho_{\delta_\star}\) 的右侧邻域内，将总和分解为
\[
\mathcal{I}_D(\varrho)
=m(\delta_\star)J_{\varrho}(\delta_\star)+R(\varrho),
\qquad
R(\varrho):=\sum_{\delta\in\Delta\setminus\{\delta_\star\}} m(\delta)\,J_{\varrho}(\delta).
\]
由 \(\delta\mapsto \varrho_\delta\) 的单射性与 \(\varrho_\star\) 的定义，\(R\) 在 \(\varrho_\star\) 的某个邻域内为 \(C^\infty\) 函数，故有 Taylor 展开
\(
R(\varrho)=R(\varrho_\star)+R'(\varrho_\star)(\varrho-\varrho_\star)+O((\varrho-\varrho_\star)^2)
\)。
另一方面，由命题 \ref{prop:xi-single-defect-threshold-phase-transition}，
\(
J_{\varrho}(\delta_\star)
=K(\delta_\star)(\varrho-\varrho_\star)^{3/2}+o((\varrho-\varrho_\star)^{3/2})
\)（\(\varrho\downarrow \varrho_\star^+\)）。
两式相加并注意
\(
\mathcal{I}_D(\varrho_\star)=R(\varrho_\star)
\)
与
\(
\mathcal{I}_D'(\varrho_\star)=R'(\varrho_\star)
\)
（因为 \(\varrho\mapsto J_{\varrho}(\delta_\star)\) 在 \(\varrho\le \varrho_\star\) 为零且在阈值处一阶导为零），即得 \eqref{eq:xi-integrated-defect-cusp-expansion}，并推出 \(A_\star=m(\delta_\star)K(\delta_\star)\)。
最后由 \(\varrho_\star=\frac{1-\delta_\star}{1+\delta_\star}\) 反解得 \(\delta_\star=\frac{1-\varrho_\star}{1+\varrho_\star}\)，且 \(m(\delta_\star)\in\NN\) 为定义中的计数。证毕。
\end{proof}

\begin{remark}[端点零化的序列判据]\label{rem:xi-integrated-defect-sequence-criterion}
由于 \(\varrho\mapsto\mathcal{I}_D(\varrho)\) 在 \((0,1)\) 上单调不减，故
\[
\lim_{\varrho\uparrow 1}\mathcal{I}_D(\varrho)=0
\quad\Longleftrightarrow\quad
\exists\,\varrho_n\uparrow 1\ \text{s.t.}\ \mathcal{I}_D(\varrho_n)\to 0.
\]
\end{remark}

\begin{corollary}[单半径平均证书给出的显式零点排除带]\label{cor:xi-integrated-defect-single-radius-exclusion-band}
\leanverified{paper\_xi\_integrated\_defect\_single\_radius\_exclusion\_band}
取任意 \(\Delta>0\) 与 \(0<\varrho<1\) 满足 \(\varrho>\varrho_{\Delta}=\frac{1-\Delta}{1+\Delta}\)（从而 \(J_{\varrho}(\Delta)>0\)）。
若
\[
\mathcal{I}_D(\varrho)\ <\ J_{\varrho}(\Delta),
\]
则不存在非平凡零点 \(\rho=\beta+\mathrm{i}\gamma\) 满足 \(\beta\le \tfrac12-\Delta\)（等价地 \(\delta(\rho)=\tfrac12-\beta\ge \Delta\)）。
\end{corollary}
\begin{proof}
反设存在 \(\rho\in\mathcal{Z}^-\) 使 \(\delta(\rho)\ge \Delta\)。
由 \eqref{eq:xi-integrated-defect-sumrule}，有
\[
\mathcal{I}_D(\varrho)\ \ge\ J_{\varrho}(\delta(\rho)).
\]
而对每个固定 \(x\in\RR\)，由 \eqref{eq:xi-cayley-modulus-delta} 可知
\(\delta\mapsto|\mathcal{C}(x+\mathrm{i}\delta)|\) 严格下降，
从而 \(\delta\mapsto \mathcal{D}_{\varrho,\delta}(x)=\bigl[\log(\varrho/|\mathcal{C}(x+\mathrm{i}\delta)|)\bigr]_+\) 单调不减；积分得 \(\delta\mapsto J_{\varrho}(\delta)\) 单调不减。
故 \(J_{\varrho}(\delta(\rho))\ge J_{\varrho}(\Delta)\)，从而 \(\mathcal{I}_D(\varrho)\ge J_{\varrho}(\Delta)\)，与假设矛盾。证毕。
\end{proof}

\begin{conjecture}[平均缺陷到一致缺陷的自举不等式]\label{conj:xi-integrated-defect-to-uniform-bootstrap}
存在绝对常数 \(A>0\) 与 \(\varrho_0\in(0,1)\)，使得对所有 \(\varrho\in(\varrho_0,1)\) 成立
\[
\sup_{x_0\in\RR}D_{\zeta,x_0}(\varrho)\ \le\ A\,(1-\varrho^2)^{-1}\,\mathcal{I}_D(\varrho).
\]
\end{conjecture}

\endinput
